\SourceLeaf{379}%
\needspace{12\baselineskip}
\begin{center}[3.]\end{center}

Die allgemeinste Lösung der Aufgabe, alle Gestalten zu bestimmen, die eine gegebene krumme Fläche annehmen kann, beruht auf Folgendem. Es sei $P$ der Punkt der Himmelskugel, dem bei einer Gestalt die Normale auf der Oberfläche der Fläche entspricht, und $Q$ der Punkt, der demselben Flächenpunkt in einer anderen Gestalt entspricht. Ebenso entsprechen die Punkte $P',P'',P'''$ usw. für eine Gestalt den Punkten $Q',Q'',Q'''$ [usw.] für die andere. Die Lösung besteht dann darin, dass jede geschlossene Figur $PP'P''P'''\, .\, .\, .$ auf der Kugeloberfläche denselben Flächeninhalt hat wie die entsprechende Figur $QQ'Q''Q'''\, .\, .\, .$.

\textls[120]{Spätere Anmerkung.} Dieses schöne Theorem ist zwar vollkommen richtig, reicht aber für die Lösung nicht aus.

\needspace{8\baselineskip}
\begin{center}BEMERKUNGEN.\end{center}

Die Notizen [2] und [3] finden sich in dem Handbuch »Den astronomischen Wissenschaften gewidmet«, das \textsc{Gauss} im November 1801 begonnen hat. Sie stammen vermutlich ungefähr aus derselben Zeit wie der unter [1] abgedruckte Brief von \textsc{Gauss} an \textsc{Schumacher}.
\begin{flushright}\textsc{Stäckel}.\end{flushright}
